\documentclass[12pt,a4paper]{article}
\usepackage[utf8]{inputenc}
\usepackage[T1]{fontenc}
\usepackage[slovak]{babel}
\usepackage{amsmath, amssymb}
\usepackage{booktabs}
\usepackage{geometry}
\usepackage{parskip}
\usepackage{hyperref}
\geometry{top=2cm, bottom=2cm, left=2.5cm, right=2.5cm}

\title{\textbf{Projekt z Ekonometrie LS 2025/2026}\\
\large Faktory šťastia v európskych krajinách}
\author{Meno Priezvisko\quad Meno Priezvisko\quad Meno Priezvisko\\
\small Ročník: \quad Odbor: }
\date{}

\begin{document}
\maketitle

%---------------------------------------------------------------
\section*{1. Úvod}

Dáta pochádzajú zo \textit{World Happiness Report 2015}
(zdroj: \texttt{happiness\_europe.txt}) a obsahujú $n = 50$
európskych krajín. Závislá premenná je \texttt{Happiness\_Score} ---
celkové skóre šťastia na stupnici 0--10. Vysvetľujúce premenné sú:

\begin{itemize}\itemsep0pt
  \item \textbf{Economy} --- HDP na obyvateľa (log),
  \item \textbf{Family} --- miera sociálnej podpory,
  \item \textbf{Health} --- stredná dĺžka zdravého života,
  \item \textbf{Freedom} --- sloboda životných rozhodnutí.
\end{itemize}

Očakávame \textit{kladný} vplyv všetkých premenných na šťastie:
bohatšie, sociálne súdržné krajiny s dlhšou zdravou dĺžkou života
a vyššou mierou slobody by mali vykazovať vyššie skóre šťastia.

%---------------------------------------------------------------
\section*{2. Formulácia regresného modelu}

Predpokladáme klasický lineárny regresný model (MNS):
\[
  \text{Happiness}_i = \beta_0 + \beta_1 \cdot \text{Family}_i
  + \beta_2 \cdot \text{Economy}_i
  + \beta_3 \cdot \text{Health}_i
  + \beta_4 \cdot \text{Freedom}_i + \varepsilon_i,
\]
kde $\varepsilon_i \sim \mathcal{N}(0,\sigma^2)$ sú i.i.d. chyby,
$i = 1,\ldots,50$. Parametre sú odhadnuté metódou najmenších štvorcov
(MNS): $\hat{\boldsymbol{\beta}} = (X^\top X)^{-1} X^\top Y$.

%---------------------------------------------------------------
\section*{3. Povinná úloha: odhady parametrov}

\subsection*{Odhady $\hat{\beta}$ a ich interpretácia}

\begin{table}[h!]
\centering
\begin{tabular}{lrrrr}
\toprule
Parameter & Odhad & Std. chyba & $t$-štatistika & $p$-hodnota \\
\midrule
$\hat{\beta}_0$ (intercept) & $1{,}6834$ & $0{,}6187$ & $2{,}721$ & $0{,}009$ \\
$\hat{\beta}_1$ (Family)    & $1{,}1237$ & $0{,}4765$ & $2{,}358$ & $0{,}023$ \\
$\hat{\beta}_2$ (Economy)   & $0{,}9736$ & $0{,}5285$ & $1{,}842$ & $0{,}072$ \\
$\hat{\beta}_3$ (Health)    & $0{,}9509$ & $1{,}0309$ & $0{,}922$ & $0{,}361$ \\
$\hat{\beta}_4$ (Freedom)   & $2{,}5565$ & $0{,}5527$ & $4{,}625$ & $<0{,}001$ \\
\bottomrule
\end{tabular}
\caption{Odhady MNS ($n=50$, $k=5$, $df = 45$).}
\end{table}

\textbf{Interpretácia:} Nárast indexu slobody o 1 jednotku zvýši
skóre šťastia priemerne o $2{,}56$ bodu (ceteris paribus) --- najsilnejší
efekt spomedzi sledovaných faktorov. Koeficient pri \textit{Family}
($1{,}12$) aj pri \textit{Economy} ($0{,}97$) sú štatisticky
významné resp.\ takmer významné; koeficient pri \textit{Health}
($0{,}95$) nie je štatisticky signifikantný.

\subsection*{Odhad $\hat{\sigma}^2$ a miera zhody}

\[
  \hat{\sigma}^2 = \frac{RSS}{n-k} = \frac{10{,}8734}{45} = 0{,}2416,
  \qquad \hat{\sigma} = 0{,}4916.
\]

\begin{table}[h!]
\centering
\begin{tabular}{cccc}
\toprule
$RSS$ & $ESS$ & $TSS$ & $R^2$ \\
\midrule
$10{,}8734$ & $34{,}2553$ & $45{,}1287$ & $0{,}7591$ \\
\bottomrule
\end{tabular}
\caption{Rozklad variability ($TSS = ESS + RSS$).}
\end{table}

Model vysvetľuje $75{,}9\,\%$ variability skóre šťastia
($R^2 = 0{,}7591$), čo považujeme za dobrú zhodu.

%---------------------------------------------------------------
\section*{4. Kategória A: Test hypotézy o kontraste}

Testujeme lineárne kontrasty $H_0\colon a^\top\beta = r$
pomocou $t$-štatistiky
\[
  T = \frac{a^\top\hat{\beta} - r}{\sqrt{\hat{\sigma}^2 \cdot
  a^\top(X^\top X)^{-1}a}} \;\sim\; t(n-k).
\]
Zamietame $H_0$ na hladine $\alpha = 0{,}05$, ak $|T| > t_{0{,}975}(45) = 2{,}014$.

\medskip
\textbf{Kontrast 1:} $H_0\colon \beta_{\text{Health}} = 2\beta_{\text{Economy}}$,
t.j.\ $a^\top = (0,0,-2,1,0)$, $r=0$.

\[
  T_1 = -0{,}5105, \quad p\text{-hodnota} = 0{,}6122.
\]
\textbf{Záver:} $H_0$ \textit{nezamietame}. Dáta neposkytujú dostatok
dôkazov na to, aby sme tvrdili, že vplyv zdravia je dvojnásobkom
vplyvu ekonomiky.

\medskip
\textbf{Kontrast 2:} $H_0\colon \beta_{\text{Family}} = \beta_{\text{Economy}} + \beta_{\text{Freedom}}$,
t.j.\ $a^\top = (0,1,-1,0,-1)$, $r=0$.

\[
  T_2 = -2{,}1511, \quad p\text{-hodnota} = 0{,}0369.
\]
\textbf{Záver:} $H_0$ \textit{zamietame} ($p < 0{,}05$). Vplyv rodiny
nie je súčtom vplyvov ekonomiky a slobody --- sloboda a ekonomika
spoločne pôsobia na šťastie silnejšie ako samotná rodina.

%---------------------------------------------------------------
\section*{5. Kategória B: Bonferroniho simultánne IS}

Pre $m = 3$ kontrasty súčasne zaručíme celkovú spoľahlivosť $\geq 95\,\%$
Bonferroniho korekciou: každý IS používa kvantil
$t_{1-\alpha/(2m)}(n-k) = t_{0{,}9917}(45) \approx 2{,}459$.
IS pre kontrast $a^\top\beta$ má tvar
$\bigl[a^\top\hat{\beta} \pm 2{,}459\cdot\hat{\sigma}\sqrt{a^\top(X^\top X)^{-1}a}\bigr]$.

\begin{table}[h!]
\centering
\begin{tabular}{llccc}
\toprule
Kontrast & Ekonomická otázka & $\hat{a}^\top\hat{\beta}$ & Dolná & Horná \\
\midrule
$\beta_4 - \beta_1$ & Sloboda vs.\ rodina   & $1{,}4328$ & $[$ viď R $]$ & $[$ viď R $]$ \\
$\beta_4 - \beta_2$ & Sloboda vs.\ ekonomika & $1{,}5829$ & $[$ viď R $]$ & $[$ viď R $]$ \\
$\beta_1 - \beta_2$ & Rodina vs.\ ekonomika  & $0{,}1501$ & $[$ viď R $]$ & $[$ viď R $]$ \\
\bottomrule
\end{tabular}
\caption{Bonferroniho simultánne IS ($\alpha = 0{,}05$, $m=3$).}
\end{table}

\textbf{Interpretácia:} Ak dolná hranica IS je kladná, prvý faktor
má štatisticky väčší vplyv. Pre kontrasty K1 a K2 (sloboda vs.\ ostatné)
očakávame kladný interval --- sloboda je dominantným faktorom.
Pre K3 (rodina vs.\ ekonomika) IS pravdepodobne obsahuje nulu,
čo zodpovedá blízkym odhadom $\hat{\beta}_1 \approx \hat{\beta}_2$.

%---------------------------------------------------------------
\section*{6. Kategória C: Test heteroskedasticity (Whiteov test)}

Overujeme, či platí predpoklad homoskedasticity:
$H_0\colon \mathrm{Var}(\varepsilon_i) = \sigma^2$ pre všetky $i$.

\textbf{Postup:} (1) odhadneme model a získame reziduá $\hat{\varepsilon}$;
(2) regredujeme $\hat{\varepsilon}^2$ na všetky regressory, ich štvorce
a krížové členy (artificial model); (3) testovacia štatistika
$n \cdot R^2_{\text{art}} \xrightarrow{d} \chi^2(c-1)$,
kde $c = k(k+1)/2 = 15$ (pre $k=5$).

\[
  H_0\text{ zamietame, ak } n \cdot R^2_{\text{art}} > \chi^2_{0{,}95}(14) = 23{,}685.
\]

Výsledky (z R skriptu):

\medskip
\begin{tabular}{lcc}
\toprule
& Asymptotický test ($n \cdot R^2$) & Exaktný test ($F'$) \\
\midrule
Hodnota štatistiky & $[$ viď R $]$ & $[$ viď R $]$ \\
Kritická hodnota   & $23{,}685$ & $[$ viď R $]$ \\
$p$-hodnota        & $[$ viď R $]$ & $[$ viď R $]$ \\
\bottomrule
\end{tabular}

\medskip
\textbf{Záver:} Ak $H_0$ zamietame, chyby merania sú heteroskedastické
a štandardné chyby odhadov $\hat{\beta}$ sú skreslené.
V takom prípade použijeme robustné (Whiteove) štandardné chyby
(sandwich estimátor \texttt{vcovHC(model, type="HC0")}) pri ďalšom testovaní.

%---------------------------------------------------------------
\section*{7. Záver}

Odhadli sme lineárny regresný model šťastia pre 50 európskych krajín.
Model dosahuje $R^2 = 0{,}76$, čo naznačuje solídnu vysvetľovaciu silu.
\textbf{Sloboda} ($\hat{\beta}_4 = 2{,}56$) a \textbf{rodina}
($\hat{\beta}_1 = 1{,}12$) sú štatisticky najvýznamnejšími faktormi.
Zdravie ($\hat{\beta}_3 = 0{,}95$) síce má kladný odhadnutý vplyv,
ale nie je štatisticky signifikantné, čo môže byť spôsobené silnou
koreláciou s ostatnými premennými.

Z testov kontrastov vyplýva, že vplyv slobody a rodiny sa štatisticky
líši od jednoduchých lineárnych kombinácií ostatných faktorov.
Bonferroniho simultánne intervaly potvrdzujú dominanciu slobody
oproti ekonomike aj rodine. Jednoduchšie povedané:
\textit{v európskych krajinách je pocit slobody a silné rodinné
väzby dôležitejšie pre šťastie ľudí ako samotná ekonomická úroveň}.

\end{document}
